% !PRINTED 339
\weberpage{339}
% running head: §. 91. | Invarianten der Curve dritter Ordnung. | 339

Diese Gleichung erhält eine einfachere Gestalt durch die
Substitution
\weq{11}{z = \frac{3m^2}{r^2}.}
Sie wird dann
\weq{12}{z^4+6Sz^2+Tz-3S^2=0.}
Diese biquadratische Gleichung hat die Eigenschaft, dass
ihre erste Invariante $A=0$ ist, und dass ihre Discriminante
(Bd. I, \S.\,47):
\weq{13}{D_2 = -27(T^2+64S^3)^2 = -27D^2,}
also stets negativ ist. Die Grösse $D=T^2+64S^3$ erhält für
die canonische Form (5) den Ausdruck:
\weq{14}{h^3(h^3+8m^3)^3,}
und kann also nicht verschwinden, so lange die Curve $f$ keinen
singulären Punkt hat. (Sie ist in der That nichts Anderes, als
die Discriminante der Form $f$.)

Aus (13) ergiebt sich nun, dass die biquadratische Gleichung
 (12) bei reellen Curven $f$ eine negative Discriminante, und
folglich (Bd. I, \S.\,78) zwei reelle und zwei conjugirt imaginäre
Wurzeln hat, und weil das Product aller vier Wurzeln $-3S^2$
negativ ist, \gesperrt{so ist von den reellen Wurzeln eine positiv},
eine negativ.

Aus $z$ kann man mit Hülfe von (11) und der ersten Gleichung
 (7) die beiden Grössen $m$ und $r$ berechnen. Man erhält so
\weq{15}{m^3 = \frac{z^2}{9S+z^2}, \quad mr^2 = \frac{3z}{9S+z^2}, \quad r^6 = \frac{27z}{(9S+z^2)^2},}
und hier könnte, wie man aus (12) schliesst, wenn man
$z^2=-9S$ setzt, $9S+z^2$ nur dann verschwinden, wenn entweder
 $S$ oder $D$ Null ist. Damit sind die Co\"efficienten $P$, $Q$, $R^3$
der cubischen Gleichung \S.\,90, (14) eindeutig durch $z$ und durch
bekannte Grössen bestimmt, und für die reelle positive Wurzel $z$
wird $m$ und $r$ reell, und folglich auch $y_1$, $y_2$, $y_3$ reell.

Eine leichte Ausnahme bildet der Fall $S=0$, in dem eine
Wurzel der Gleichung (12) verschwindet, und die zweite reelle
Wurzel $-\sqrt[3]{T}$ wird, und für die zugehörigen Werthe von $m$
und $r^6$ erhält man aus (7) und (11)
\[
m = 0, 1; \quad r^6 = \frac{1}{T}, \ \frac{-27}{T};
\]
%CONT>>
